\documentclass[11pt]{article}

\usepackage[margin=1in]{geometry}
\usepackage{booktabs}
\usepackage{amsmath}
\usepackage{hyperref}
\usepackage{xcolor}
\usepackage{array}
\usepackage{longtable}
\usepackage{parskip}
\usepackage{caption}
\usepackage{fancyhdr}
\usepackage{titlesec}
\usepackage[htt]{hyphenat}

\titleformat{\section}{\normalfont\Large\bfseries}{\thesection}{1em}{}
\titleformat{\subsection}{\normalfont\large\bfseries}{\thesubsection}{1em}{}

\pagestyle{fancy}
\fancyhf{}
\fancyhead[L]{Prediction Accuracy Report v3}
\fancyhead[R]{\thepage}
\renewcommand{\headrulewidth}{0.4pt}

\hypersetup{
    colorlinks=true,
    linkcolor=blue,
    urlcolor=blue,
    citecolor=blue,
    pdftitle={Prediction Algorithm Accuracy Report v3 - Multi-Season Results},
    pdfauthor={Sports Analytics Project}
}

\newcommand{\code}[1]{\allowbreak\texttt{\small #1}\allowbreak}

\title{\textbf{Prediction Algorithm Accuracy Report — v3} \\ \large Multi-Season Data: What More Data Actually Bought Us}
\author{Sports Analytics Project}
\date{September 2026}

\begin{document}

\maketitle

\begin{abstract}
\noindent This report covers the third round of work on the platform's prediction algorithms, following two
single-season reports (\code{prediction-accuracy-report.pdf}, \code{-v2.pdf}, same directory) that concluded
recency-weighted feature recombination was exhausted and the two highest-value remaining bets were a genuine
data-correctness fix and more historical data. Both were executed: a real trade-attribution bug was found and
fixed (players' historical box scores were being attributed to their \emph{current} team, not the team they
played for in that specific game), and two additional complete NBA seasons were ingested, taking the dataset
from 1,321 games to 3,781. Re-running the Elo constant grid search against the enlarged dataset produced a
different, validated result from the single-season attempt: tuned constants
(\code{HOME\_COURT\_ADVANTAGE\_ELO}\,=\,40, \code{K\_FACTOR}\,=\,25) now beat the literature defaults on
held-out data, confirmed stable across a wide parameter grid and two independent validation splits, and have
been shipped. Every other model improved as well, most strikingly the two recency-weighted player-point
predictors, whose edge over a naive baseline roughly doubled. We also document two engineering incidents from
the ingestion process itself — a Postgres connection pooler killing idle/long-lived connections far more
aggressively than expected — since the fixes are reusable for any future data pull.
\end{abstract}

\tableofcontents

\section{What Changed Since v2}

\begin{table}[h]
\centering
\small
\renewcommand{\arraystretch}{1.3}
\begin{tabular}{@{}>{\raggedright\arraybackslash}p{0.30\textwidth}>{\raggedright\arraybackslash}p{0.13\textwidth}>{\raggedright\arraybackslash}p{0.51\textwidth}@{}}
\toprule
Change & Status & Detail \\
\midrule
Trade-attribution bug fix & \textbf{Shipped} & New \code{PlayerGameStat.teamId} column, schema migration, ingestion fix, backfill of 33,590 existing rows, \code{four\_factors.py} query fix (\S\ref{sec:trade-bug}) \\
2023-24 season ingested & \textbf{Shipped} & 1,230 games, full boxscores, correct team attribution from day one \\
2024-25 season ingested & \textbf{Shipped} & 1,230 games, same \\
Elo constants re-tuned & \textbf{Shipped} & \code{HOME\_COURT\_ADVANTAGE\_ELO} 75\,\textrightarrow\,40, \code{K\_FACTOR} 20\,\textrightarrow\,25 — this time validated, unlike the v1/v2 attempt (\S\ref{sec:elo-retune}) \\
Ingestion connection resilience & \textbf{Shipped} & \code{ingest\_historical\_season.py} — fresh DB connection per game write, not one held across the whole run (\S\ref{sec:pooler}) \\
\bottomrule
\end{tabular}
\caption{Disposition of this round's changes.}
\end{table}

\section{The Trade-Attribution Bug}
\label{sec:trade-bug}

While scoping the multi-season pull, we found that \code{four\_factors.py}'s team-game boxscore aggregation
joined \code{PlayerGameStat} to \code{Team} via \code{Player.teamId} — a player's \emph{current} team — rather
than the team they actually suited up for in a specific historical game. A player traded mid-season had every
one of their pre-trade games silently reattributed to their new team the moment the roster ingestion next ran.
An audit query confirmed this was not theoretical: \textbf{7.7\% of the 33,590 \code{PlayerGameStat} rows}
already had a \code{Player.teamId} that matched neither team in that row's own game.

The fix required a real schema change, not just a query edit, since the correct information (which team each
boxscore row's player was on, for that specific game) was being read from the NBA API and then discarded
before ever reaching the database:

\begin{enumerate}
    \item Added \code{PlayerGameStat.teamId} via a Prisma migration (nullable, foreign-keyed to \code{Team}).
    \item Updated \code{apps/ingestion/games.py}/\code{ingest.py} to capture and write this field going
    forward — the boxscore endpoint already reports it per player, it just wasn't being kept.
    \item Backfilled all 33,590 existing rows: 31,012 were unambiguous (the player's current team matched one
    of the two teams in that game, so no fix was needed) and were updated directly; the remaining 2,578 rows
    (774 distinct games) needed a boxscore re-fetch to resolve correctly, since the original ingestion never
    stored per-game team identity.
    \item Updated \code{four\_factors.py}'s query to join on the new field.
\end{enumerate}

\textbf{Measured impact}, isolated from the data-volume increase (re-run against the original single-season
dataset before any new games were added): MAE improved from 12.52 to 12.51 points, walk-forward MAE from
12.65 to 12.62. A small, real improvement — expected, since only 7.7\% of rows were affected and the fix
redistributes existing signal between two teams' season averages rather than adding new information. The fix
was worth making on correctness grounds regardless of its accuracy impact, and matters increasingly as more
seasons (with more real trades) are added.

\section{The Multi-Season Elo Re-Tuning: A Different Result This Time}
\label{sec:elo-retune}

\subsection{What changed}

v1/v2 reported that a chronological train/validation grid search over \code{HOME\_COURT\_ADVANTAGE\_ELO} and
\code{K\_FACTOR}, run against the single-season dataset (1,321 games), picked a training-split winner that
validated \emph{worse} than the FiveThirtyEight literature defaults (HCA=75, K=20) — textbook overfitting, not
a real improvement, attributed to too little independent evidence in one season to beat a well-chosen prior.

With the dataset now at 3,781 games (2023-24, 2024-25, 2025-26), the identical grid search was re-run.

\begin{table}[h]
\centering
\begin{tabular}{lrr}
\toprule
Split & Default (HCA=75, K=20) Brier & Tuned (HCA=40, K=25) Brier \\
\midrule
70/30 chronological split & 0.2142 & 0.2119 \\
60/40 chronological split & 0.2123 & 0.2101 \\
\bottomrule
\end{tabular}
\caption{Validation-split Brier score, default vs.\ tuned constants, multi-season dataset. Lower is better;
tuned constants win on both independent splits this time.}
\end{table}

Two checks were run before trusting this, given the exact same experiment produced the opposite conclusion
on less data:

\begin{itemize}
    \item \textbf{Grid stability.} The winning pair (HCA=40, K=25) was re-confirmed against a much wider grid
    (HCA $\in [0, 110]$, K $\in [5, 45]$) and sits comfortably in the interior, not pinned against either
    boundary — a boundary-pinned winner would suggest the true optimum lies outside the tested range and the
    result shouldn't be trusted as-is.
    \item \textbf{Split independence.} The same winner was reproduced from two different chronological
    train/validation cut points (70/30 and 60/40 of the season timeline), not just the one split used in the
    original single-season report. Both splits favor the tuned constants over the defaults by a similar
    margin, which is the standard cross-validation argument that this is a real effect, not an artifact of
    exactly where one particular split happened to fall.
\end{itemize}

The training-split Brier surface itself also tightened meaningfully with more data — the top-10 grid
candidates spanned 0.2233–0.2251 on the single-season dataset (a difference small enough to be mostly noise)
versus 0.2156–0.2162 on the 3-season dataset (a narrower, more confidently-identified optimum). This is
exactly the signature you'd expect if the earlier result really was a data-volume problem rather than a
methodology problem, as the v1/v2 report speculated.

\subsection{What shipped}

\code{HOME\_COURT\_ADVANTAGE\_ELO} was changed from 75 to 40, and \code{K\_FACTOR} from 20 to 25, in
\code{elo.py}. \code{GamePrediction} was regenerated for all 3,781 games under the new constants.

\begin{table}[h]
\centering
\begin{tabular}{lrr}
\toprule
Metric & Single-season (old constants) & Multi-season (tuned constants) \\
\midrule
$n$ & 1,321 & 3,781 \\
Brier score & 0.2146 & 0.2145 \\
Accuracy (threshold 0.5) & 65.2\% & 66.0\% \\
Naive home-favorite baseline & 55.5\% & 54.8\% \\
\bottomrule
\end{tabular}
\caption{Full-dataset accuracy, old vs.\ new. Brier score is nearly unchanged on the full (non-held-out)
figure — the real gain is in calibration (Table~\ref{tab:calibration}) and in how the model generalizes to
unseen games (the validation-split numbers above), not in a large full-dataset summary-statistic move.}
\end{table}

\begin{table}[h]
\centering
\small
\begin{tabular}{lrrrr}
\toprule
Predicted range & $n$ (old) & Old: pred / actual & $n$ (new) & New: pred / actual \\
\midrule
0.2--0.3 & 27 & 0.257 / 0.037 & 311 & 0.253 / 0.296 \\
0.3--0.4 & 87 & 0.357 / 0.230 & 421 & 0.352 / 0.363 \\
0.4--0.5 & 181 & 0.454 / 0.359 & 572 & 0.453 / 0.442 \\
\bottomrule
\end{tabular}
\caption{Calibration in the range where v1 found the worst overconfidence. The predicted-vs-actual gap
(e.g.\ predicting 0.454 when the true rate was 0.359, a 0.095 overconfidence) has closed substantially at the
new constants (0.453 predicted vs.\ 0.442 actual, a 0.011 gap) — this is the practically meaningful
improvement, even though the single headline Brier number barely moved.}
\label{tab:calibration}
\end{table}

The permutation test (1,000 randomized-outcome reruns of the real schedule) again rejects the null decisively
($p < 0.001$) — the model's accuracy is real signal, not an artifact of the larger dataset alone.

\subsection{A caveat we deliberately did not address}

\code{compute\_elo\_ratings} treats every game across all three seasons as one continuous chronological
sequence — a team's rating carries directly from the last game of one season into the first game of the next,
with no reset or regression-toward-the-mean for the roster turnover a real NBA offseason involves (trades,
draft, free agency). This is a real simplification. It was deliberately left unaddressed in this round so the
re-tuning experiment isolated one variable (more data under the existing model) rather than two (more data
\emph{and} a season-boundary model change) — conflating them would have left it unclear which change actually
produced the improvement. A season-aware reset is a reasonable candidate for a future, separately-evaluated
change.

\section{Impact on the Other Three Models}

\begin{table}[h]
\centering
\small
\begin{tabular}{lrr}
\toprule
Model & Single-season & Multi-season \\
\midrule
Four Factors margin MAE & 12.52 & 12.30 \\
Four Factors naive-baseline MAE & 13.36 & 12.91 \\
Fantasy points: recency-weighted MAE & 7.74 & 8.00 \\
Fantasy points: naive-mean MAE & 8.25 & 8.94 \\
Fantasy points: recency-weighting edge & +0.52 & +0.94 \\
Scorer points: recency-weighted MAE & 4.45 & 4.63 \\
Scorer points: naive-mean MAE & 4.64 & 4.98 \\
Scorer points: recency-weighting edge & +0.20 & +0.36 \\
\bottomrule
\end{tabular}
\caption{All four models, single- vs.\ multi-season. MAE in points (or fantasy points) throughout.}
\end{table}

\textbf{Four Factors} improved modestly (12.52 $\to$ 12.30 MAE) — a real gain, consistent with more team-game
observations sharpening the same 3-feature regression, but not a transformation the way Elo's calibration
was. It remains the weakest of the four models relative to its naive baseline and should stay a secondary,
lower-confidence signal in the product.

\textbf{Both recency-weighted player-point predictors got measurably better relative to their own naive
baseline}, even though their raw MAE went up slightly. This is not a contradiction: MAE going up while the
\emph{edge over naive} roughly doubles means the naive baseline got substantially worse with more data (a
longer unweighted history dilutes toward a player's career-long average, increasingly stale as more games
accumulate), while the recency-weighted model held up better against the same harder task. In other words,
recency-weighting's specific value proposition — tracking current form instead of career-long average —
becomes more important, not less, exactly as more historical data is added. This is the strongest positive
result in this report for a fully-shipped, no-caveats model.

\section{Engineering Note: Pooled-Connection Fragility During Ingestion}
\label{sec:pooler}

Two separate incidents during the historical-season ingestion are worth recording since the underlying cause
recurs for any future long-running script against this project's Supabase database.

\textbf{Incident 1 — idle connection killed during a slow, mostly-network-bound phase.} The original
\code{ingest\_historical\_season.py} opened one database connection at the start of the script and held it
through \code{collect\_season\_game\_dates}, a phase making 30 rate-limited \code{LeagueGameFinder} calls with
no database activity in between. That connection was killed by Supabase's pooler before a single write
happened, with the error \texttt{server closed the connection unexpectedly}.

\textbf{Incident 2 — connection killed mid-write-loop, even while actively used.} After fixing incident 1 by
opening a fresh connection only for the write phase, the run still failed, repeatedly, partway through —
including within the first 50 games. A direct test isolated the cause precisely: \textbf{a connection idle
for as little as 5 seconds was already killed} (\texttt{cur.execute('SELECT 1')} on a 5-second-idle connection
failed with the same error). Each game in the write loop involves a rate-limited $\sim$1 second-plus network
call to \code{stats.nba.com} \emph{before} the next database write — comfortably enough idle time, every
single game, to exceed a 5-second pooler timeout.

\textbf{Fix.} Rather than reconnect-on-failure (tried first; the pooler was frequent enough that even a
retry-once policy failed multiple times within the first 50 games), the script now opens, writes, commits,
and closes a fresh connection for \emph{every individual game}. This adds one extra network round-trip per
game but means every write is durably committed the instant it happens — a crash or interruption loses at
most one in-flight game, and simply re-running the script (every write is an idempotent upsert) picks up
exactly where it left off. This pattern generalizes to any future long-running script against this database:
do not hold one connection across a loop with rate-limited network calls between writes; treat this pooler's
practical idle budget as a few seconds, not the tens of seconds or minutes that might be assumed for a
typical managed Postgres instance.

\section{Recommendations Going Forward}

\begin{enumerate}
    \item \textbf{Treat the Elo re-tuning result as a positive signal for pulling even more data}, specifically
    for Four Factors and the two player-point predictors, none of which improved as dramatically as Elo did.
    Elo's constants were the most data-starved parameter in the whole system (only 2 free parameters,
    literature priors already close), so it's the one most likely to show a clean signal early; the others
    may simply need more than 3 seasons to show a similarly clear multi-season effect.
    \item \textbf{Revisit the season-boundary reset question} (\S\ref{sec:elo-retune}'s caveat) as a focused,
    separately-evaluated follow-up — now that Elo's core constants are validated against real multi-season
    data, adding a principled offseason regression-toward-the-mean is the next natural refinement, and can be
    backtested the same way (train/validation split, compare against the current no-reset baseline).
    \item \textbf{Apply the same trade-attribution correctness discipline going forward}: any future feature
    that joins boxscore data to a team should use \code{PlayerGameStat.teamId}, not \code{Player.teamId} —
    worth a code-review note or lint rule if more code ends up needing this join.
    \item \textbf{Reuse the fresh-connection-per-write pattern} (\S\ref{sec:pooler}) for any future ingestion
    or long-running write script against this database, rather than rediscovering the 5-second idle timeout
    the hard way again.
\end{enumerate}

\section*{Appendix: Reproducing These Results}

\code{apps/predictor/check\_accuracy.py} recomputes the Elo and Four Factors numbers in this report against
the database's current state — run it directly rather than trusting the point-in-time figures above once
more data has been added. The player-point predictor backtests (recency-weighting edge over a naive running
mean) and the Elo grid-search/validation-split methodology are one-off analysis scripts, not part of the
committed codebase, reflecting their purpose as a one-time validation exercise rather than ongoing monitoring
— \code{check\_accuracy.py} is the durable, re-runnable piece.

\end{document}
